\subsection{Lastenheft}
\label{app:Lastenheft}

Das Lastenheft beschreibt die fachlichen Anforderungen des Auftraggebers an das
\Fachbegriff{AntispamAdminCenter}. Ziel ist die zentrale Verwaltung von
Firmen, Benutzern, Mail-Adressen und IP-Adressen f"ur die
Antispam-nahe Mail-Infrastruktur.

\subsubsection*{Ausgangssituation}
Vor Projektbeginn wurden Freigaben und Zust"andigkeiten nicht in einer
zentralen Anwendung gepflegt. Relevante Informationen waren auf mehrere
Datenquellen verteilt. Dadurch entstanden Mehrfacheingaben, Medienbr"uche,
Abstimmungsaufwand und eine nur eingeschr"ankt nachvollziehbare Historie von
"Anderungen.

\subsubsection*{Zielsetzung}
Die neue Anwendung soll die Pflege vereinheitlichen, die
Mandantentrennung sicher abbilden und fachliche "Anderungen revisionsf"ahig
protokollieren. Die L"osung soll intern "uber den Browser nutzbar sein und eine
zuverl"assige Datenbasis f"ur nachgelagerte Mail-Systeme bereitstellen.

\subsubsection*{Funktionale Anforderungen}
\begin{enumerate}
\item Die Anwendung muss Firmen zentral anlegen, "andern, anzeigen und
      verwalten k"onnen.
\item Die Anwendung muss Benutzer mit Rollen und Firmenzuordnung
      verwalten k"onnen.
\item Die Anwendung muss Mail-Adressen firmenbezogen erfassen, validieren
      und bearbeiten k"onnen.
\item Die Anwendung muss IP-Adressen firmenbezogen erfassen, validieren
      und bearbeiten k"onnen.
\item Die Anwendung muss zwischen globalen und mandantenbezogenen Rechten
      unterscheiden.
\item Die Anwendung muss eine abgesicherte Anmeldung f"ur berechtigte Benutzer
      bereitstellen.
\item Die Anwendung muss fachliche "Anderungen an den verwalteten Objekten in
      einem Audit-Log protokollieren.
\item Die Anwendung muss die verwalteten Mail- und IP-Daten f"ur nachgelagerte
      technische Verarbeitung strukturiert bereitstellen.
\item Die Anwendung soll eine einheitliche und leicht verst"andliche
      Bedienoberfl"ache f"ur die t"agliche Pflege bieten.
\end{enumerate}

\subsubsection*{Nichtfunktionale Anforderungen}
\begin{enumerate}
\item Die Anwendung muss ohne lokale Client-Installation im internen Netzwerk
      "uber einen Webbrowser erreichbar sein.
\item Die Anwendung muss eine klare Trennung von Datenmodell, Fachlogik und
      Weboberfl"ache besitzen, damit sie wartbar erweitert werden kann.
\item Die Anwendung muss fehlerhafte Eingaben fr"uhzeitig erkennen und dem
      Benutzer verst"andlich zur"uckmelden.
\item Die Anwendung muss sicherheitsrelevante Zugriffsvorg"ange angemessen
      absichern, insbesondere durch Passwort-Hashing, Session-Verwaltung und
      Schutz gegen wiederholte Fehlanmeldungen.
\item Die Anwendung soll f"ur den vorgesehenen internen Benutzerkreis mit
      vertretbarem Administrationsaufwand betrieben werden k"onnen.
\end{enumerate}

\subsubsection*{Abgrenzung}
Im Projekt wird nicht der Spamfilter selbst entwickelt. Es geht um die
Verwaltung der daf"ur ben"otigten Stamm- und Freigabedaten. Auch der endg"ultige
Produktivbetrieb, eine umfangreiche "Ubernahme alter Daten aus Fremdsystemen
und zus"atzliche Spezial-Schnittstellen zu Mail-Systemen geh"oren nicht zum
Projektumfang.
